\section{Proyección con los radios de cada sitio}
Sea $c$ el número de componentes del grafo de incidencias de $U$.
Se usa el rango cíclico del grafo bipartito de sitios y fibras ocupadas, con las
fibras unitarias como hojas. Definimos $p=v+c+1,v+c,v+c-1$ según el número de
ciclos independientes sea cero, uno o mayor que uno.
Al proyectar en todas las activaciones, el coste es a lo sumo $2^v$.
Si un campo se anula, las dos caras de activación de ese sitio coinciden.
Así se puede aplicar además $\prod_i(I-\operatorname{ev}_{\zeta_i=0})$.
Un holomorfo acotado por $M$ en el polidisco de radios $\kappa'/d_i$,
evaluado tras esta proyección con $|\zeta_i|\le C/j$, cuesta a lo sumo
\[
 M\prod_i\frac{Cd_i/j}{1-Cd_i/j}\le M\prod_i(C' d_i/j).
\]
Sumamos los monomios en que todos los sitios aparecen. Si al menos uno se repite,
se paga además $\sum_i C d_i/j\le Cv^2/j$ por unión sobre el sitio repetido.
El radio en la variable de tamaño es $\kappa/v$. Los órdenes a partir de $w$ en $s$
cuestan a lo sumo $2(Cv/j)^w$ con $s=1/j$, sobre el mismo polidisco de campos.
Las cancelaciones algebraicas de \emph{Grados intermedios}, parte 05, se aplican
a los mismos aproximantes, ahora estimados en discos mayores. El bosque empieza en
$w=c+1$. Con varios ciclos empieza en $w=c-1$. Con un ciclo los órdenes ordinarios
empiezan en $w=c$, y el jet $c-1$ tiene una actividad repetida como mínimo.
Ese jet cuesta $(Cv/j)^{c-1}Cv^2/j$ además de los $v$ campos. Para el error
físico se usan sólo las actividades; $\Delta_q\le C j^{-q/2}$ y $q/2\ge c+1$.
Resulta en los tres casos
\[
       |[U]V|\le C^v (Cv)^{c+1}\left(\prod_i d_i\right) j^{-p},
        \qquad 5\le v\le L\log n .
\]
El $O(v)$ de tamaño del holomorfo cabe en $C^v$. Si $U$ es colineal, la
sustracción de la línea da cero. Si existe un sitio sin vecino diagonal, su
activación es indiferente y el coeficiente se anula. Las sustracciones de parejas
tienen soporte a lo sumo cuatro. Los grados menores conservan su cota archivada.
